\documentstyle[% 


righttag% 


,11pt% 


,amssymb% 


,russian%               remove in Eng. 


,izvru%                 Set izven% in Eng. 


]{amsart} 


 


% 


 


\newcommand{\tts}[1]{} 


 


 


%\hoffset=-15mm%                  For Eng. 


\setcounter{page}{61} 


\god{2004} 


\nomer{4~(503)} %                        Remove in Eng. 


%\nomer{Vol.\,48, No.\,4}%               Set in Eng. 


%\str{\pageref{begin}--\pageref{end}}%  Set in Eng. 


\udk{517.956} 


 


\author{\it Е.А.\,Уткина} 


% NB:   ^^ remove in Eng. 


\title{О задачах Гурса с дополнительными нормальными производными 


в краевых условиях} 


 


\date{} 


%\pagestyle{myheadings}%         Set in Eng. 


\pagestyle{plain} %              Remove in Eng. 


\begin{document} 


%\thispagestyle{plain}%          Set in Eng. 


\maketitle 


\label{begin} 


 


 


 


Пусть $D$ --- параллелепипед $x_0<x<x_1$, $y_0<y<y_1$, $z_0<z<z_1$, а 


$X$, $Y$, $Z$ --- его грани при $x=x_0$, $y=y_0$, $z=z_0$ соответственно. 


Будем говорить 


об одном трехмерном аналоге уравнений из [1] 


\begin{equation} 


L(u)\equiv\sum^2_{\alpha,\beta=0}\sum^1_{\gamma=0} 


a_{\alpha\beta\gamma}(x,y,z) 


\frac{\partial^{\alpha+\beta+\gamma}u} 


{\partial x^\alpha\partial y^\beta\partial z^\gamma}=0. 


\end{equation} 


В работе [2] для (1) решена задача Гурса об отыскании в $D$ функции 


$u(x,y,z)$ по условиям 


\begin{gather} 


u|_X=\varphi(y,z),\quad u|_Y=\psi(x,z),\quad u|_Z=\theta(x,y),\notag\\ 


u_x|_X=\varphi_1(y,z),\quad u_y|_Y=\psi_1(x,z). 


\end{gather} 


 


В данной статье рассматриваются задачи, которые соотносятся с 


(1), (2) аналогично тому, как в теории эллиптических уравнений задача 


Неймана соотносится с задачей Дирихле. А именно, 


эти задачи получаются заменой в 


(2) хотя бы одного граничного условия значением нормальной производной из 


набора 


\begin{equation} 


u_{xx}(x_0,y,z)=\varphi_2(y,z),\quad 


u_{yy}(x,y_0,z)=\psi_2(x,z),\quad 


u_z(x,y,z_0)=\theta_1(x,y). 


\end{equation} 


 


Как и в [1], считаем, что $a_{221}(x,y,z)\equiv1$, а условия гладкости 


остальных 


коэффициентов (1) определяются включениями $a_{\alpha\beta\gamma}\in 


C^{\alpha+\beta+\gamma}(\ov D)$, 


где $C^{\alpha+\beta+\gamma}$ есть класс непрерывных в $\ov D$ вместе с их 


производными 


$\partial^{r+s+t}/\partial x^r\partial y^s \partial z^t$ ($r=0,\dotsc,\alpha$; 


$s=0,\dotsc,\beta$; $t=0,\dotsc,\gamma$) функций. Указанные задачи 


можно считать обобщением ситуаций, исследованных для более 


простых уравнений в работах [2], [3]. Отметим еще, что уравнения из [1], 


[3], принадлежащие этому же классу, встречаются в приложениях (уравнения 


Адлера и Буссинеска--Лява). 


 


Присутствие условия из (3) на соответствующей характеристике 


будем обозначать буквой $N$, а отсутствие --- $\varGamma$. Получаемые при 


этом 


варианты естественно закодировать сочетанием указанных букв, например, 


$\varGamma\varGamma N$, $NNN$ и т.\,д. 


 


Остановимся на трех задачах. 


\medskip 


 


{\bf Задача $N\varGamma\varGamma$.} {\it Найти функцию $u\in C^{2+2+1}(D) 


\cap C^{2+0+0}(D\cup\ov X) \cap C^{0+1+0}(D\cup\ov Y) \cap 


C^{0+0+0}(D\cup\ov Z)$, 


являющуюся в $D$ решением уравнения $(1)$, 


удовлетворяющую условию 


$$ 


u_{xx}(x_0,y,z)=\varphi_2(y,z),\quad 


\varphi_2\in C^{2+1}(\ov X),


$$ 


и всем соотношениям $(2)$, кроме первого.} 


\medskip 


 


Решение может быть основано на редукции к задаче (1), (2), осуществляемой 


следующим образом. Дважды проинтегрируем (1) по $y$ в пределах 


от $y_*$ до $y$ ($(x,y_*,z),(x,y,z)\in D$), затем в полученном соотношении 


устремим $y_*$ к $y_0$. Далее интегрируем один раз по $z$ в пределах от 


$z_*$ до $z$ и устремляем $z_*$ к $z_0$, а $x$ к $x_0$. Учитывая граничные 


условия, 


получим 


\begin{equation} 


\sum^2_{i=0}\sum^1_{j=0}D^{-i}_y D^{-j}_z 


(A_{0ij}\varphi(y,z))=r_1(y,z), 


\end{equation} 


\begin{gather*} 


r_1(y,z)=-\sum^2_{i=0}\sum^1_{j=0}D^{-i}_y D^{-j}_z 


(A_{2ij}\varphi_2(y,z)+A_{1ij}\varphi_1(y,z))+ \\ 


% 


+\sum^2_{i=0}D^{-i}_y  


(A_{2i0}\varphi_2(y,z_0)+A_{1i0}\varphi_2(y,z_0)+ 


A_{0i0}\theta(x_0,y))+ \\ 


% 


+\sum^1_{i=0}\sum^1_{j=0}D^{-i}_y D^{-j}_z


(A_{2ij}\varphi_2(y_0,z)+A_{1ij}\varphi_1(y_0,z)+ 


A_{0ij}\psi(x_0,z))- \\ 


% 


-\sum^1_{i=0}D^{-i}_y  


(A_{2i0}\varphi_2(y_0,z_0)+A_{1i0}\varphi_1(y_0,z_0)+ 


A_{0i0}\theta(x_0,y_0))+ \\ 


% 


+D^{-1}_y(\varphi_{2y}(y_0,z)-\varphi_{2y}(y_0,z_0))+ 


D^{-1}_y D^{-1}_z(a_{220}(x_0,y_0,z) 


\varphi_{2y}(y_0,z))+ \\ 


% 


+D^{-2}_y D^{-1}_z(a_{121}(x_0,y_0,z)\varphi_{1y}(y_0,z))+ 


D^{-1}_y(a_{121}(x_0,y_0,z_0)\varphi_{1y}(y_0,z_0)- \\ 


% 


-a_{121}(x_0,y_0,z_0)\varphi_{1y}(y_0,z_0))- \\ 


% 


-D^{-1}_y D^{-1}_z(a_{121z}(x_0,y_0,z)\varphi_{1y}(y_0,z)- 


a_{120}(x_0,y_0,z)\varphi_{1y}(y_0,z))+ \\ 


% 


+D^{-1}_y(a_{021}(x_0,y_0,z)\psi_1(x_0,z)- 


a_{021}(x_0,y_0,z_0)\psi_1(x_0,z_0))- \\ 


% 


-D^{-1}_y D^{-1}_z(a_{021z}(x_0,y_0,z)\psi_1(x_0,z)- 


a_{020}(x_0,y_0,z)\psi_1(x_0,z)). 


\end{gather*} 


Здесь $D^k_t\varphi\equiv\partial^k\varphi/\partial t^k$ при 


$k=1,2,\dotsc$ и $D^k_t\varphi\equiv \int\limits^t_{t_0} 


\frac{(t-\tau)^{-k-1}\varphi(\tau)}{(-k-1)!}d\tau$ при


$k=-1,-2,\dotsc$, $D^0_t$ --- единичный оператор; 


 


$A_{i00}(y,z)=a_{i21}(x_0,y,z)$, 


$A_{i10}(y,z)=-2a_{i21y}(x_0,y,z)+a_{i11}(x_0,y,z)$, 


 


$A_{i20}(y,z)=a_{i21yy}(x_0,y,z)-a_{i11y}(x_0,y,z)+ a_{i00}(x_0,y,z)$, 


 


$A_{i01}(y,z)=-a_{i21z}(x_0,y,z)+a_{i20}(x_0,y,z)$, 


 


$A_{i11}(y,z)=-a_{i01z}(x_0,y,z)+a_{i10}(x_0,y,z)- 


2(-a_{i21yz}(x_0,y,z)+a_{i20y}(x_0,y,z))$, 


 


$A_{i21}(y,z)=-a_{i21yz}(x_0,y,z)+a_{i20y}(x_0,y,z)+ 


a_{i11yz}(x_0,y,z)-a_{i10y}(x_0,y,z)- a_{i01z}(x_0,y,z)+a_{i00}(x_0,y,z)$. 


 


При этом если $y=y_0$, то там, где было умножение на 2, совершается 


умножение на 1; $\varphi(y,z)$ --- функция из первого условия (2), которую 


в данном 


случае требуется определить. 


 


Из уравнения (4) на основании рассуждений, аналогичных приведенным 


в [2], [3], нетрудно усмотреть, что имеет место 


 


\begin{thm} 


Выполнение неравенства 


\begin{equation} 


\sum^2_{j=0}\sum^1_{k=0}A^2_{0jk}\ne0 


\end{equation} 


обеспечивает однозначное определение функции $\varphi(y,z)$ через 


$\varphi_2(y,z)$, 


$\varphi_1(y,z)$, $\psi(x,z)$, $\psi_1(x,z)$, $\theta(x,y)$. При этом 


функция $\varphi(y,z)$ записывается 


в явном виде в любом из следующих случаев\/${:}$ 


 


$1)$ обе функции $A_{0jk}$, $A_{0j+1k}$ отличны от нуля при $j=0,1$, 


$k=0,1;$ 


 


$2)$ обе функции $A_{0jk}$, $A_{0jk+1}$ отличны от нуля при $j=0,1,2$, 


$k=0;$ 


 


$3)$ $A_{00k}\ne0$, $A_{00k}-y A_{02k}\equiv0$, $A_{02k}\ne0$ при $k=0,1;$ 


 


$4)$ $A_{0jk}$ $(j=0,1,2;$ $k=0,1)$ в нуль не обращается, все остальные 


коэффициенты равны нулю тождественно. 


\end{thm} 


 


Включение $\varphi\in C^{2+1}(\ov X)$ обеспечивается добавлением к 


условиям 


гладкости на коэффициенты (1) требований $a_{i_1j_1k_1}\in 


C^{0+(j_1+j)+(k_1+k)}(D\cup\ov X)$. 


 


Приведем вариант расшифровки указанных выше условий. Например, 


выполнение 1) при $j=0$, $k=0$ означает, что $a_{021}(x_0,y,z)\ne0$, 


$A_{010}(y,z)\ne0$, $A_{020}(y,z)\equiv0$, $A_{001}(y,z)\equiv0$, 


$A_{011}(y,z)=-a_{011z}(x_0,y,z)+a_{010}(x_0,y,z)\equiv0$, 


$A_{021}(y,z)=-a_{001z}(x_0,y,z)+a_{000}(x_0,y,z)\equiv0$. 


Явный вид $\varphi$ задается в этом случае формулой 


$$ 


\varphi(y,z)=\frac{1}{A_{000}(y,z)}\bigg\{r_1(y,z) 


-\int^y_{y_0}A_{010}(\eta,z)r_1(\eta,z)\bigg[\exp 


\int^\eta_y A_{010}(\tau,z)d\tau\bigg]d\eta\bigg\}. 


$$ 


 


 


 


{\bf Задача $\varGamma N\varGamma$.} {\it Отличается от предыдущей тем, 


что в $(2)$ 


второе условие заменяется на 


$$ 


u_{yy}(x,y_0,z)=\psi_2(x,z),\quad 


\psi_2\in C^{2+1}(\ov Y). 


$$ 


Здесь $u\in C^{2+2+1}(D) \cap C^{1+0+0}(D\cup\ov X) \cap 


C^{0+2+0}(D\cup\ov Y)\cap C^{0+0+0}(D\cup\ov Z)$.} 


\medskip 


 


 


Так как переменные $x$ и $y$ входят в (1) симметрично, для 


нахождения $\psi(x,z)$ получим уравнение, подобное (4). Вместо $A_{ijk}$ 


здесь будут 


использоваться $B_{jik}$, полученные перестановкой первых двух индексов. 


 


Анализ этого уравнения показывает, что справедлива 


 


\begin{thm} 


Достаточным условием однозначной разрешимости 


задачи $\varGamma N\varGamma$ является неравенство $(5)$ с заменой 


$A_{0jk}$ на $B_{j0k}$. При этом 


редукция к задаче Гурса обеспечивается выполнением условий $1)$--\,$4)$ из 


формулировки теоремы $1$, в которых $A_{0jk}$ заменены на $B_{j0k}$. 


\end{thm} 


 


Условия гладкости на коэффициенты в данном случае определяются 


включениями


$$


a_{i_1j_1k_1}\in C^{(i_1+i)+0+(k_1+k)}(D\cup\ov Y). 


$$


 


{\bf Задача $\varGamma\varGamma N$.} {\it Отыскать $u\in C^{2+2+1}(D) \cap 


C^{1+0+0}(D\cup\ov X)\cap C^{0+1+0}(D\cup\ov Y)\cap C^{0+0+1}(D\cup\ov 


Z)$, 


являющуюся решением $(1)$ при выполнении всех 


условий $(2)$, кроме третьего, вместо которого берется} 


$$ 


u_z(x,y,z_0)=\theta_1(x,y),\quad \theta_1\in C^{2+2}(\ov Z). 


$$ 


 


 


Здесь при выводе аналога (4) интегрируем (1) дважды по $x$ и 


дважды по $y$. Тогда для определения $\theta(x,y)$ получим уравнение 


\begin{equation} 


\sum^2_{i=0}\sum^2_{j=0}D^{-i}_x D^{-j}_y 


(C_{ij0}\theta(x,y))=r_3(x,y), 


\end{equation} 


где 


\begin{gather*} 


r_3(x,y)=-\sum^2_{i=0}\sum^2_{j=0}D^{-i}_x D^{-j}_y 


(C_{ij1}\theta_1(x,y))+\sum^2_{i=0}\sum^1_{j=0} 


D^{-i}_x D^{-j}_y(C_{ij0}\psi(x,z_0)+ 


C_{ij1}\theta_1(x,y_0))+ \\ 


% 


+\sum^1_{i=0}\sum^2_{j=0}D^{-i}_x D^{-j}_y 


(C_{ij0}\varphi(y,z_0)+C_{ij1}\theta_1(x_0,y))- 


D^{-1}_x D^{-1}_y(\theta_{1xy}(x_0,y_0))- \\ 


% 


-\sum^1_{i=0}\sum^1_{j=0}D^{-i}_x D^{-j}_y 


(C_{ij0}\varphi(y_0,z_0)+C_{ij1}\theta_1(y_0,z_0))+ \\ 


% 


+D^{-1}_y(\psi_{1z}(x,z_0)-\psi_{1z}(x_0,z_0))+ 


D^{-1}_x(\theta_{1x}(x_0,y)-\theta_{1x}(x_0,y_0))+ \\ 


% 


+D^{-1}_y(a_{220}(x,y_0,z_0)\psi_1(x,z_0)- 


a_{220}(x_0,y_0,z_0)\psi_1(x_0,z_0))+ \\ 


% 


\displaybreak[3]


+D^{-1}_x(a_{220}(x_0,y,z_0)\varphi_1(y,z_0)- 


a_{220}(x_0,y_0,z_0)\varphi_1(y_0,z_0))- \\ 


% 


-D^{-1}_x D^{-1}_y(2a_{220x}(x,y_0,z_0)\psi_1(x,z_0)- 


a_{220x}(x_0,y_0,z_0)\psi_1(x_0,z_0)+ \\ 


% 


+2a_{220y}(x_0,y,z_0)\varphi_1(y,z_0)- 


a_{220x}(x_0,y_0,z_0)\psi_1(x_0,z_0)- \\ 


% 


-a_{220y}(x_0,y_0,z_0)\varphi_1(y_0,z_0)+ 


a_{211}(x_0,y,z_0)\varphi_{1z}(y,z_0)- 


a_{211}(x_0,y_0,z_0)\varphi_{1z}(y_0,z_0))+ \\ 


% 


+D^{-1}_x D^{-2}_y((a_{220yy}(x_0,y,z_0)- 


a_{210y}(x_0,y,z_0)+a_{200}(x_0,y,z_0))\varphi_1(y,z_0)+ \\ 


% 


+a_{201yy}(x_0,y,z_0)\varphi_{1z}(y,z_0)+ 


a_{200}(x_0,y,z_0)\varphi_1(y,z_0))+ \\ 


% 


+D^{-2}_x D^{-1}_y((a_{220xx}(x,y_0,z_0)- 


a_{120x}(x,y_0,z_0)+a_{020x}(x,y_0,z_0))\psi_1(x,z_0)+ \\ 


% 


+(-a_{121x}(x,y_0,z_0)+a_{021}(x,y_0,z_0))\psi_{1z}(x,z_0))+ \\ 


% 


+D^{-1}_x D^{-1}_y(a_{210}(x_0,y,z_0)\varphi_1(y,z_0)- 


a_{210}(x_0,y_0,z_0)\varphi_1(y_0,z_0)+ \\ 


% 


+ a_{121}(x,y_0,z_0)\psi_{1z}(x,z_0)- 


a_{121}(x_0,y_0,z_0)\psi_{1z}(x_0,z_0)+ \\ 


% 


+a_{120}(x,y_0,z_0)\psi_1(x,z_0)- 


a_{120}(x_0,y_0,z_0)\psi_1(y_0,z_0)), 


\end{gather*} 


а $C_{kls}$ задаются формулами 


\begin{gather*} 


C_{00i}(x,y)=a_{20i}(x,y,z_0),\quad 


C_{01i}=-2a_{22iy}(x,y,z_0)+a_{21i}(x,y,z_0), \\ 


% 


C_{02i}=a_{22iyy}(x,y,z_0)-a_{21iy}(x,y,z_0)+ 


a_{20i}(x,y,z_0), \\ 


% 


C_{10i}=-2a_{22ix}(x,y,z_0)+a_{12i}(x,y,z_0),\quad 


C_{20i}=a_{22ixx}(x,y,z_0)-a_{12ix}(x,y,z_0)+ 


a_{02i}(x,y,z_0),\\ 


% 


C_{11i}=4a_{22ixy}(x,y,z_0)-2a_{21ix}(x,y,z_0)- 


2a_{12iy}(x,y,z_0)+a_{11i}(x,y,z_0),\\ 


% 


C_{21i}=-2a_{22ixxy}(x,y,z_0)+a_{21ixx}(x,y,z_0)+ \\+


2a_{12ixy}(x,y,z_0)-a_{11ix}(x,y,z_0)- 


2a_{02iy}(x,y,z_0)+a_{01i}(x,y,z_0), \\ 


% 


C_{12i}=-2a_{22ixyy}(x,y,z_0)+2a_{21ixy}(x,y,z_0)+ \\+


a_{12iyy}(x,y,z_0)-a_{11iy}(x,y,z_0)-2a_{20ix}(x,y,z_0) 


+a_{10i}(x,y,z_0), \\ 


% 


C_{22i}=-a_{12ixyy}(x,y,z_0)+a_{22ixxyy}(x,y,z_0)+ 


a_{11ixy}(x,y,z_0)-a_{10ix}(x,y,z_0)+ \\ 


% 


+a_{20ixx}(x,y,z_0)-a_{21ixxy}(x,y,z_0)-a_{01iy}(x,y,z_0)+ 


a_{02iyy}(x,y,z_0)+a_{00i}(x,y,z_0). 


\end{gather*} 


При этом если $x=x_0$ и была производная по $x$, то в тех слагаемых, где 


производилось умножение на 2 или на 4, соответственно выполняется 


умножение на 1 или 2. Аналогично для $y=y_0$. В случае обоих равенств 


все коэффициенты равны 1. 


 


Аналогично предыдущим случаям выводится 


 


\begin{thm} 


Функция $\theta(x,y)$ однозначно определяется через 


$\theta_1(x,y)$, $\varphi(y,z)$, $\varphi_1(y,z)$, $\psi(x,z)$, 


$\psi_1(x,z)$ при выполнении условия 


$$ 


\sum^2_{i=0}\sum^2_{j=0}C^2_{ij0}\ne0. 


$$ 


 


При этом $\theta(x,y)$ записывается в явном виде в любом из следующих 


случаев\/${:}$ 


\begin{itemize} 


\item[1)] обе функции $C_{ij0}$, $C_{i+1j0}$ отличны от нуля при $i=0,1$, 


$j=0,1,2;$ 


\item[2)] обе функции $C_{ij0}$, $C_{ij+10}$ отличны от нуля при 


$i=0,1,2$, $j=0,1;$ 


\item[3)] $C_{0j0}\ne0$, $C_{0j0}-xC_{2j0}\equiv0$, $C_{2j0}\ne0$ при 


$j=0,1,2$, $C_{i00}\ne0$, $C_{i00}-yC_{i20}\equiv0$, $C_{i20}\ne0$ при 


$i=0,1,2;$ 


\item[4)] только одна функция $C_{ij0}$ $(i=0,1,2;$ $j=0,1,2)$ отлична от 


нуля, а все остальные 


равны нулю. Включение $\theta\in C^{2+2}(\ov Z)$ обеспечивается 


добавлением 


к уже имеющимся в $[2]$ условиям гладкости на коэффициенты $(1)$ 


требований $a_{i_1j_1k_1}\in C^{(i_1+i)+(j_1+j)+0}(D\cup\ov Z)$. 


\end{itemize} 


\end{thm} 


 


Рассмотренные задачи можно считать базовыми: разрешимость 


остальных указанных в начале статьи легко получается с их помощью. 


Поясним это на одной из них. 


\medskip 


 


{\bf Задача $NNN$.} {\it Найти для $(1)$ решение класса 


$u\in C^{2+2+1}(D) \cap C^{2+0+0}(D\cup\ov X) \cap C^{0+2+0}(D\cup\ov Y) 


\cap C^{0+0+1}(D\cup\ov Z)$, 


удовлетворяющее условиям $(3)$ и двум последним условиям в $(2)$.} 


\medskip 


 


Функции $\varphi(y,z)$, $\psi(x,z)$, $\theta(x,y)$ определяются из 


уравнений (4), 


(6). При этом в каждом случае для нахождения неизвестной функции 


используются неравенства типа (5). В общем случае редукция к задаче Гурса 


осуществляется в терминах резольвент. Для нахождения решения в явном 


виде требуется комбинировать варианты 1)--4) теорем 1--3 на $X$, $Y$, $Z$ 


с 


учетом различных $i$, $j$, $k$ подобно тому, как это делается в [3]. Всего 


комбинаций будет $4^3$: 111, 121, 211, $112,\dotsc,444$ (пишем номера подряд 


без скобок, на первом месте указан соответствующий пункт из теоремы~1, 


на втором --- из теоремы 2, на третьем --- из теоремы 3). Каждый набор 


дает 


информацию об условиях гладкости и числе входящих в решение произвольных 


функций и (или) констант. Число различных произвольных 


функций в решении может доходить до трех. 


 


\begin{thebibliography}{9} 


 


\bibitem{1} 


Солдатов А.П., Шхануков М.Х. {\em Краевые задачи с общим нелокальным 


условием А.А.\,Самарского для псевдопараболических уравнений 


высокого порядка} // 


ДАН СССР. -- 1987. -- Т.\,297. -- \No\,3. -- С.\,547--552. 


\bibitem{2} 


Жегалов В.И., Уткина Е.А. {\em Задача Гурса для одного трехмерного 


уравнения 


со старшей производной} // Изв. вузов. Математика. 


-- 2001. -- \No\,11. -- С.\,77--81. 


\bibitem{3} 


Жегалов В.И., Миронов А.Н. {\em Трехмерные характеристические задачи с 


нормальными производными в граничных условиях} // 


Дифференц. уравнения. -- 2000. -- 


Т.\,36. -- \No\,6. -- С.\,833--836. 


\bibitem{4} 


Жегалов В.И., Уткина Е.А. {\em Об одном псевдопараболическом уравнении 


третьего порядк}а // Изв. вузов. Математика. -- 1999. -- \No\,10. 


-- С.\,73--76. 


 


\end{thebibliography} 


 


\vskip 2cm 


 


\post{\it Казанский государственный}{\it Поступила} 


\post{\it педагогический университет}{05.12.2001} 


 


\label{end} 


\end{document} 


